\section{Case Study: When Retrieval Decides the Answer}
\label{app:case-nausea}
\setlength{\emergencystretch}{2em}

This appendix records the full transcripts behind Figure~\ref{fig:case-nausea}
(\S\ref{sec:retr-corpus}): two near-identical queries whose answers diverge
because retrieval handed the generator different chunks. The system's
\emph{queries} and \emph{generated answers} are reproduced verbatim, as the
deployed system's own input and output. The retrieved chunks are reproduced as
the \emph{actual text the model received} (lightly cleaned of PDF-extraction
artifacts); we withhold all source titles and page numbers and instead tag the
four distinct source documents \mbox{``A''--``D''} (full provenance is in the
corpus manifest). \textbf{The labels are per source document}: the same letter
means the same document, possibly a different passage from it. That re-use is
the crux---source~A supplies both the management protocol Case~A retrieves and
the bare definition Case~B retrieves, and source~D supplies two further
passages---so the same guidance can be surfaced or missed depending only on
which passage the retriever returns. Bracketed numbers \mbox{[1]--[3]} mark the
retrieval order and match the citations in the generated answers.

\subsection*{Case A --- query says ``severe nausea''}

\noindent\textbf{Query (verbatim).} \emph{``A 20-year-old woman, Gravida~1,
currently in her second trimester of pregnancy with an estimated due date of
15 September, is suffering from severe nausea. What are the recommended
treatment options?''}

\smallskip
\noindent\textbf{Retrieved top-3 (actual chunk text; sources withheld).}

\smallskip
\noindent\emph{[1] Source~A --- management section.}
\begin{quote}\footnotesize
\textbf{Moderate vomiting and hyperemesis}

If the woman is experiencing \emph{moderate} vomiting or suspected
\emph{hyperemesis gravidarum}, the midwife should arrange for her to be admitted
to hospital without delay. She should be cared for in a single room if possible,
to avoid undue disturbance. A urinalysis is conducted for ketones, bilirubin,
protein and glucose, and a midstream specimen of urine is sent for culture to
exclude pyelonephritis.

Antihistamines such as \emph{promethazine hydrochloride, prochlorperazine
(Stemetil) metoclopramide (Maxolon, Primperan)} (Jewell \& Young 2003) or
corticosteroids such as \emph{methylprednisolone} (Ismail \& Kenny 2007) are
considered safe to use in pregnancy and may be utilized to control the nausea
and vomiting.

Blood tests, such as aspartate aminotransferase (AST), alanine aminotransferase
(ALT), urea and electrolytes, are required to assess renal and liver function.
An ultrasound scan may also be useful to eliminate hydatidiform mole and
multiple pregnancy as probable causes. A frequent record of the woman's weight,
temperature, pulse and blood pressure is essential to monitor her wellbeing. In
addition, the fetal heart rate should be appropriately auscultated (depending on
the gestation) to monitor the health of the fetus.
\end{quote}

\noindent\emph{[2] Source~D --- rationale.}
\begin{quote}\footnotesize
\textbf{Why the committee made the recommendations}

gravidarum can have serious harmful consequences, and treatment and care in
hospital may be needed. It should be noted that this guideline only covers
treatments to manage nausea and vomiting in pregnancy, and comprehensive
management of hyperemesis gravidarum, which may include nutritional
interventions, is not covered by this guideline on routine antenatal care.
\end{quote}

\noindent\emph{[3] Source~D --- recommendations.}
\begin{quote}\footnotesize
\textbf{Nausea and vomiting}
\begin{itemize}\setlength{\itemsep}{1pt}
  \item 1.4.1 Reassure women that mild to moderate nausea and vomiting are
  common in pregnancy, and are likely to resolve before 16 to 20 weeks.
  \item 1.4.2 Recognise that by the time women seek advice from healthcare
  professionals about nausea and vomiting in pregnancy, they may have already
  tried a number of different interventions.
  \item 1.4.3 For pregnant women with mild-to-moderate nausea and vomiting who
  prefer a non-pharmacological option, suggest that they try ginger.
\end{itemize}
\end{quote}

\smallskip
\noindent\textbf{Generated answer (verbatim).}
\begin{quote}\small
Severe nausea in pregnancy needs careful management.
\begin{itemize}\setlength{\itemsep}{0pt}
  \item \textbf{Mild to moderate nausea:} Reassure the woman that this is
  common. Suggest trying \textbf{ginger} as a non-drug option~[3].
  \item \textbf{Severe nausea (Hyperemesis Gravidarum):} If the nausea is
  severe, she needs hospital care without delay~[1].
\end{itemize}
\emph{Hospital Care Steps:}
\begin{enumerate}\setlength{\itemsep}{0pt}
  \item Monitor her \textbf{weight, temperature, pulse, and blood pressure}
  frequently~[1].
  \item Check her \textbf{fetal heart rate} regularly~[1].
  \item Perform a \textbf{urinalysis} for ketones, bilirubin, protein, and
  glucose~[1].
  \item Send a \textbf{midstream urine specimen} for culture to check for
  infection~[1].
  \item Blood tests for \textbf{liver and kidney function} (AST, ALT, urea,
  electrolytes) are needed~[1].
  \item An \textbf{ultrasound scan} may be used to rule out other causes~[1].
\end{enumerate}
\emph{Medication Options:}
\begin{itemize}\setlength{\itemsep}{0pt}
  \item Antihistamines like \textbf{promethazine hydrochloride,
  prochlorperazine (Stemetil), metoclopramide (Maxolon, Primperan)}, or
  corticosteroids like \textbf{methylprednisolone} may be used to control
  symptoms~[1].
\end{itemize}
Please consult a doctor or senior clinician for the final treatment plan.
\end{quote}

\subsection*{Case B --- query says ``hyperemesis gravidarum''}

\noindent\textbf{Query (verbatim).} \emph{``A 20-year-old woman, Gravida~1,
currently in her second trimester of pregnancy with an estimated due date of
15 September, is suffering from hyperemesis gravidarum. What are the
recommended treatment options?''} (The same patient and condition as Case~A;
only the wording differs.)

\smallskip
\noindent\textbf{Retrieved top-3 (actual chunk text; sources withheld).}

\smallskip
\noindent\emph{[1] Source~B --- general management section.}
\begin{quote}\footnotesize
\textbf{Management}

The woman is usually admitted to hospital for investigations and rehydration.
The cause of vomiting may not be found. An \textbf{antiemetic} will be given and
fluids and electrolytes replaced by intravenous infusion of a solution such as
Hartmann's. Vitamin B12, thiamine, vitamin C, folic acid and iron will be needed
to prevent complications (Farrell 2008).

General observation of the woman's condition should be monitored. Strict fluid
balance should be maintained until rehydrated. There is usually a rapid response
to treatment and oral fluids may be recommenced when vomiting has ceased for
24~h. Solid food should be then introduced gradually. Moran \& Taylor (2002)
found that weight loss of more than 5\% of prepregnancy weight in women with HG
was effectively treated with 10~mg of prednisolone three times a day. This
shortened the stay in hospital and stopped vomiting. This treatment was
gradually decreased and discontinued at 20 weeks of gestation.
\end{quote}

\noindent\emph{[2] Source~C --- student study-exercise.}
\begin{quote}\footnotesize
\textbf{Activity}

Find out what is meant by hyperemesis gravidarum.

See if you can discover what is the incidence and who is most at risk of this
condition.

Find out how the midwife would recognize this condition, and what the treatment
is.
\end{quote}

\noindent\emph{[3] Source~A --- definition section} (the same document as
Case~A's passage~[1]).
\begin{quote}\footnotesize
\textbf{Hyperemesis gravidarum}

\emph{Hyperemesis gravidarum} is a pathological condition characterized by
unremitting, severe vomiting in pregnancy. It occurs in 0.3--2.0\% of
pregnancies and is more common in women who are younger, non-smokers and
non-Caucasian (Ismail \& Kenny 2007). It is diagnosed by exclusion (Kametas \&
Nelson-Piercy 2008) and is a leading cause of hospital admission during
pregnancy (Cedergren et al 2008).

Studies have shown that a gastric infection caused by the \emph{Helicobacter
pylori} bacterium may also be linked to hyperemesis gravidarum (Golberg et al
2007). The infection responds well to antibiotic therapy, such as
\emph{erythromycin} or \emph{clarithromycin}.
\end{quote}

\smallskip
\noindent\textbf{Generated answer (verbatim).}
\begin{quote}\small
Hyperemesis gravidarum is severe, continuous vomiting during pregnancy. For
this woman, the recommended treatment options include:
\begin{itemize}\setlength{\itemsep}{0pt}
  \item \textbf{Hospital Admission:} She usually needs to be admitted for
  tests and rehydration.
  \item \textbf{Fluid and Electrolyte Replacement:} Give intravenous fluids,
  such as Hartmann's solution.
  \item \textbf{Medication:} An \textbf{antiemetic} will be given.
  \item \textbf{Supplements:} Give Vitamin B12, thiamine, Vitamin C, folic
  acid, and iron to prevent problems.
  \item \textbf{Monitoring:} Watch her condition closely and keep her fluid
  balance strict until she is rehydrated.
  \item \textbf{Diet:} Start oral fluids when vomiting stops for 24 hours.
  Then, slowly introduce solid food.
\end{itemize}
If she loses more than 5\% of her pre-pregnancy weight, a doctor may consider
\textbf{prednisolone}~[2]. Please consult a doctor for the exact treatment
plan and medication doses.
\end{quote}

\smallskip
\noindent\textbf{Takeaway.} The two queries describe the same patient and the
same condition; only the wording differs. Case~A surfaced the management
section and produced a specific, actionable answer with named drugs and a full
workup; Case~B surfaced the \emph{definition section} of the same source, a
generic management passage whose antiemetic is unnamed, and a study-exercise,
and the answer collapsed to ``an antiemetic will be given.'' The corpus held
the right management section in both cases---a small, faithful generator
reproduced whatever it was handed. This is the per-query face of
\S\ref{sec:retr-corpus}: corpus coverage, specificity, curation, and query
phrasing decide deployment quality.
